\renewcommand*{\authorOfNextLines}{Helmut Quinz}
\begin{quote}
	\gqqEmph{Glück ist, was passiert, wenn Vorbereitung auf Gelegenheit trifft.}\\
	\small{Seneca d.J.}
\end{quote}
Die Vorbereitung dient zur Findung und grundsätzlichen Bewertung von Projektideen.
Die Aufgabe hier sind 
\begin{itemize}
	\item Das Beschreiben der Ideen
	\item Festlegen welche Ziele mit dem Projekt \textbf{nicht} verfolgt werden
	\item Die Art des Projektes identifizieren. Solche könnten sein
		\begin{itemize}
			\item Entwickeln eines vollständig neuen Produktes oder Verfahren
			\item Kombination bereits verfügbarer Technologien zu einer neuen Funktionseinheit
			\item Optimierung einer bereits bestehenden Funktionseinheit oder eines bestehenden Verfahrens
			\item \dots
		\end{itemize} 
		Auch ist hier zu klären ob es sich 
		\begin{itemize}
			\item reines Softwareprojekt,
			\item ein Projekt zur Konzeptentwicklung,
			\item ein Simulationsprojekt oder 
			\item ein Projekt  mit Hardware-Bedarf
		\end{itemize}
		handelt
	\item Analyse des Projektumfelds. Dazu gehören 
		\begin{itemize}
			\item die Identifikation von Stakeholdern, also Personen die auf das Projekt Einfluss nehmen können oder davon betroffen sind.
					Dazu gehören auch betreuende Lehrkräfte!
			\item Analyse des Projektrisikos.
					Welche Sachverhalte oder Ereignisse können sich negativ auf die Erreichung des Projektzieles auswirken.
					Wie kann diesen entgegengewirkt werden.
			\item Ressourcenverfügbarkeit.
					Stehen alle für den Projekterfolg erforderlichen Mittel (finanziell, personell, technologisch, \dots) zu Verfügung.
					Gibt es irgendwo Engpässe oder Einschränkungen
			\item Gesetzliche und regulatorische Anforderungen.
					Ist das Projekt im Einklang mit bestehenden gesetzlichen Bestimmungen und allgemeinen Standards umsetzbar?
		\end{itemize}
	\item das Beschreiben des, sich aus einem positiven Projektergebnisses ergebenden Nutzens für den Geldgeber (Business Case)
	\item Festlegen eines Projektleiters
\end{itemize}
\paragraph{Wichtig} Sie sollten Ihre Projektidee in dieser Phase mit möglichen Betreuern diskutieren!